\documentclass[aspectratio=169,11pt]{beamer}

\usepackage[T1]{fontenc}
\usepackage[utf8]{inputenc}
\usepackage[brazil]{babel}
\usepackage{lmodern}
\usepackage{booktabs}
\usepackage{amsmath}
\usepackage{graphicx}
\usepackage{tikz}
\usetikzlibrary{positioning}

\definecolor{azulSO}{HTML}{173B57}
\definecolor{azulClaro}{HTML}{2D7FA3}
\definecolor{verdeSO}{HTML}{2A8C6A}
\definecolor{laranjaSO}{HTML}{D97A2B}
\definecolor{cinzaSO}{HTML}{F2F5F7}
\definecolor{textoSO}{HTML}{263238}

\usetheme{default}
\setbeamertemplate{navigation symbols}{}
\setbeamercolor{normal text}{fg=textoSO,bg=white}
\setbeamercolor{structure}{fg=azulSO}
\setbeamercolor{frametitle}{fg=white,bg=azulSO}
\setbeamercolor{block title}{fg=white,bg=azulClaro}
\setbeamercolor{block body}{fg=textoSO,bg=cinzaSO}
\setbeamercolor{alerted text}{fg=laranjaSO}
\setbeamerfont{frametitle}{series=\bfseries}
\setbeamertemplate{itemize item}{\color{azulClaro}$\blacktriangleright$}
\setbeamertemplate{itemize subitem}{\color{verdeSO}$\bullet$}
\setbeamertemplate{footline}{%
  \leavevmode\hbox{%
    \begin{beamercolorbox}[wd=.78\paperwidth,ht=2.7ex,dp=1.1ex,leftskip=1em]{author in head/foot}%
      \color{azulSO}\insertshorttitle
    \end{beamercolorbox}%
    \begin{beamercolorbox}[wd=.22\paperwidth,ht=2.7ex,dp=1.1ex,rightskip=1em plus1fil]{date in head/foot}%
      \color{azulSO}\hfill\insertframenumber/\inserttotalframenumber
    \end{beamercolorbox}}}

% Os PNGs ficam ao lado deste arquivo no projeto do Overleaf.
% Os limites simultâneos preservam a proporção, sem cortar ou deformar.
\newcommand{\figuraresultado}[1]{%
  \includegraphics[width=.94\textwidth,height=.68\textheight,keepaspectratio]{#1}}

\title[Simulador de Escalonamento]{Simulador de Escalonamento de Processos}
\subtitle{Sistemas Operacionais — Projeto da Unidade 3}
\author[Equipe]{Elder Rayan Oliveira Silva \quad Espedito Ramom Mascena Ricarto \\
Manoel Junio Duarte da Silva \quad Pedro Yan Alcantara Palácio \\
Sabrina Alencar Soares \quad Samuel Wagner Tiburi Silveira \\
Sebastião Sousa Soares}
\date{Agosto de 2026}

\begin{document}

\begin{frame}[plain]
  \titlepage
  \vfill
  \centering
  \small Universidade Federal do Cariri — UFCA
\end{frame}

\begin{frame}{Por que comparar escalonadores?}
  \begin{columns}[T,onlytextwidth]
    \column{.55\textwidth}
      \begin{block}{Motivação}
        A política de escalonamento afeta desempenho, justiça e custo operacional do sistema.
      \end{block}
      \begin{itemize}
        \item Simular estados, filas, E/S e preempção
        \item Comparar políticas com métricas quantitativas
        \item Avaliar um algoritmo próprio de forma reproduzível
      \end{itemize}
    \column{.42\textwidth}
      \begin{block}{Fluxo do projeto}
        \centering
        Carga com seed\\[.4em]
        $\downarrow$\\[.4em]
        Simulação em ticks\\[.4em]
        $\downarrow$\\[.4em]
        Métricas + IC95\%\\[.4em]
        $\downarrow$\\[.4em]
        Análise crítica
      \end{block}
  \end{columns}
  \vfill
  \scriptsize Troca de contexto: 1 tick entre PIDs distintos; não conta no primeiro despacho nem após ociosidade.
\end{frame}

\begin{frame}{Modelo de processo e de E/S}
  \centering
  \begin{tikzpicture}[node distance=5mm and 4mm,
    state/.style={draw=azulClaro,rounded corners,fill=cinzaSO,minimum width=19mm,minimum height=8mm,font=\small},
    arrow/.style={->,thick,color=azulSO}]
    \node[state] (n) {Novo};
    \node[state,right=of n] (p) {Pronto};
    \node[state,right=of p] (e) {Executando};
    \node[state,right=of e] (b) {Bloqueado};
    \node[state,right=of b] (f) {Finalizado};
    \draw[arrow] (n)--(p); \draw[arrow] (p)--(e);
    \draw[arrow] (e)--(b); \draw[arrow,bend left=28] (b) to (p);
    \draw[arrow] (e) to[bend left=20] (f);
  \end{tikzpicture}
  \vspace{.8em}
  \begin{columns}[T,onlytextwidth]
    \column{.49\textwidth}
      \begin{block}{Processo}
        PID, chegada, prioridade, estado e rajadas alternadas:
        \[\text{CPU}\rightarrow\text{E/S}\rightarrow\text{CPU}\]
        Fila de prontos contém apenas processos aptos.
      \end{block}
    \column{.49\textwidth}
      \begin{block}{Escolha de modelagem}
        E/S paralela entre processos, sem dispositivo ou fila própria. O retorno ocorre ao fim da duração sorteada.

        \alert{Limitação:} não mede contenção de dispositivos e pode favorecer cargas I/O-bound.
      \end{block}
  \end{columns}
\end{frame}

\begin{frame}{Cargas iguais, comparação justa}
  \begin{columns}[T,onlytextwidth]
    \column{.52\textwidth}
      \begin{block}{Reprodutibilidade}
        \begin{itemize}
          \item PCG32 controlado por seed de 64 bits
          \item Chegadas acumuladas em intervalos de 0 a 3 ticks
          \item Workload CSV reutilizado por todos os algoritmos
          \item SHA-256 identifica a carga exata
          \item Execução em lote pode ser interrompida e retomada
        \end{itemize}
      \end{block}
    \column{.45\textwidth}
      \begin{block}{Quatro cenários obrigatórios}
        \begin{enumerate}
          \item Equilibrado
          \item I/O-bound
          \item CPU-bound
          \item Prioridades desbalanceadas
        \end{enumerate}
      \end{block}
      \centering
      \textbf{seed + cenário}\quad$\Longrightarrow$\quad\textbf{mesma carga}
  \end{columns}
\end{frame}

\begin{frame}{Algoritmos de referência}
  \small
  \begin{tabular}{@{}p{.17\textwidth}p{.48\textwidth}p{.22\textwidth}@{}}
    \toprule
    \textbf{Política} & \textbf{Critério de escolha} & \textbf{Preempção} \\
    \midrule
    FCFS & Maior tempo contínuo na fila & Não \\
    Round Robin & Fila FIFO; quantum padrão de 4 ticks & Sim, por quantum \\
    Prioridade & Menor valor numérico (0 é mais prioritário) & Não \\
    SJF extra & Menor próxima rajada estimada pelo histórico & Não \\
    \bottomrule
  \end{tabular}
  \vspace{1em}
  \begin{block}{Decisões que tornam a comparação reproduzível}
    Empates usam tempo na fila e PID; todos recebem os mesmos eventos, custos e cargas. Nenhum escalonador conhece rajadas futuras, CPU restante ou processos ainda não admitidos.
  \end{block}
\end{frame}

\begin{frame}{PDBH: Prioridade Dinâmica Baseada em Histórico}
  \begin{columns}[T,onlytextwidth]
    \column{.54\textwidth}
      \begin{block}{Problema e proposta}
        Mitigar inanição da prioridade estática combinando \textit{aging} e penalidade pelo uso observado da CPU.
      \end{block}
      \[
        \text{score}=\text{prioridade}
        -\left\lfloor\frac{\text{espera}}{10}\right\rfloor
        +\left\lfloor\frac{\text{CPU consumida}}{5}\right\rfloor
      \]
      Menor score vence; empate por tempo na fila e PID.
    \column{.43\textwidth}
      \begin{block}{O que o diferencia}
        \begin{itemize}
          \item Recalcula a ordem em cada despacho
          \item Usa apenas presente e histórico
          \item Não altera a prioridade original
          \item Permanece não preemptivo
        \end{itemize}
      \end{block}
      \alert{Hipótese:} beneficiar quem espera e penalizar quem já consumiu mais CPU.
  \end{columns}
\end{frame}

\begin{frame}{Como avaliamos desempenho e justiça}
  \begin{columns}[T,onlytextwidth]
    \column{.52\textwidth}
      \begin{align*}
        \text{turnaround}_i &= C_i-A_i \\
        \text{ideal}_i &= \sum CPU_i+\sum E/S_i \\
        \text{slowdown}_i &= \frac{\text{turnaround}_i}{\text{ideal}_i}
      \end{align*}
      \begin{block}{Trocas de contexto}
        Mudanças diretas entre PIDs: custo potencial do escalonador.
      \end{block}
    \column{.45\textwidth}
      \begin{block}{Justiça de Jain}
        \[
          J=\frac{(\sum s_i)^2}{n\sum s_i^2}\times100
        \]
        Próximo de 100\% indica slowdowns semelhantes.
      \end{block}
      \alert{Justiça não basta:} todos podem ser igualmente lentos. Jain deve ser interpretado junto ao turnaround.
  \end{columns}
\end{frame}

\begin{frame}{Protocolo experimental e auditoria}
  \begin{columns}[T,onlytextwidth]
    \column{.55\textwidth}
      \centering
      {\Large\bfseries 4 cenários $\times$ 100 seeds}\par
      \vspace{.3em}
      {\Large\bfseries $\times$ 5 algoritmos = 2.000 execuções}\par
      \vspace{.8em}
      {\large 1.000 processos por carga}\par
      \vspace{.8em}
      Média e intervalo de confiança de 95\%
    \column{.42\textwidth}
      \begin{block}{Evidências reproduzíveis}
        \begin{itemize}
          \item Configuração e seeds versionadas
          \item Manifestos e hashes
          \item CSVs consolidados
          \item Testes unitários e integração
          \item Gráficos gerados por script
        \end{itemize}
      \end{block}
  \end{columns}
\end{frame}

\begin{frame}{Resultado 1: turnaround médio}
  \centering
  \figuraresultado{mean_turnaround_ic95.png}

  \small \textbf{Leitura:} Prioridade obteve o menor turnaround em todos os cenários; RR apresentou o maior, especialmente em CPU-bound. Barras mostram média e IC95\%, em ticks.
\end{frame}

\begin{frame}{Resultado 2: justiça de Jain do slowdown}
  \centering
  \figuraresultado{jain_slowdown_pct_ic95.png}

  \small \textbf{Leitura:} RR é mais justo; Prioridade amplia a desigualdade. O eixo vertical está em porcentagem e as hastes representam IC95\%.
\end{frame}

\begin{frame}{Resultado 3: trocas de contexto}
  \centering
  \figuraresultado{context_switches_ic95.png}

  \small \textbf{Leitura:} a preempção do RR eleva o custo, sobretudo em CPU-bound. O eixo vertical representa quantidade de trocas e as hastes, IC95\%.
\end{frame}

\begin{frame}{Conclusão: o PDBH respondeu à hipótese}
  \begin{columns}[T,onlytextwidth]
    \column{.55\textwidth}
      \begin{block}{O que os dados mostraram}
        \begin{itemize}
          \item PDBH não melhorou nenhuma métrica ou cenário
          \item IC95\% praticamente sobreposto ao FCFS
          \item Mesmo número de trocas das políticas não preemptivas
          \item Custo extra: recalcular a pontuação da fila
        \end{itemize}
      \end{block}
    \column{.42\textwidth}
      \begin{block}{Por quê?}
        Sem preempção, processos esperam juntos e acumulam bônus semelhantes. A pontuação tende a preservar a ordem de chegada.
      \end{block}
      \alert{Conclusão defendida:} histórico dinâmico isolado não basta. Próximo passo: investigar um PDBH preemptivo.
  \end{columns}
  \vfill
  \centering\small Código testado • resultados auditáveis • artigo científico • responsabilidades registradas
\end{frame}

\end{document}
